\documentclass[11pt,a4paper]{article}
\usepackage[utf8]{inputenc}
\usepackage[T1]{fontenc}
\usepackage[brazilian]{babel}
\usepackage{lmodern}
\usepackage{amsmath,amssymb}
\usepackage[margin=2.25cm]{geometry}
\usepackage{booktabs,array}
\usepackage{xcolor}
\usepackage{fancyhdr}
\usepackage{listings}
\usepackage{tcolorbox}
\usepackage[hidelinks]{hyperref}

\definecolor{azul}{HTML}{183B56}
\definecolor{azulClaro}{HTML}{EAF2F8}
\definecolor{cinza}{RGB}{245,245,245}
\definecolor{cinzaEscuro}{RGB}{80,80,80}
\definecolor{verdeEscuro}{HTML}{1E7E34}

\pagestyle{fancy}
\fancyhf{}
\fancyhead[L]{\small Projeto e Análise de Algoritmos -- Prof. Felipe A. Louza}
\fancyhead[R]{\small Lista 7: Seleção (Estudo e Revisão)}
\fancyfoot[L]{\small Arthur Pires Lopes -- 12221EEL001}
\fancyfoot[R]{\thepage}
\setlength{\headheight}{14pt}
\setlength{\parindent}{0pt}
\setlength{\parskip}{6pt}
\setlength{\emergencystretch}{2em}
\renewcommand{\arraystretch}{1.3}

\newcommand{\res}[1]{\[\boxed{#1}\]}
\newcommand{\topico}[1]{\par\medskip\textbf{\color{azul}#1}\par}
\newcommand{\questao}[2]{\section*{Questão #1 -- #2}}
\newcommand{\destaque}[1]{\textbf{\color{azul}#1}}

\newtcolorbox{caixaEstudo}[1][]{
  colback=azulClaro!50,
  colframe=azul,
  fonttitle=\bfseries\color{white},
  title={#1},
  arc=3mm,
  boxrule=1pt,
  left=8pt, right=8pt, top=6pt, bottom=6pt
}

\lstset{
  basicstyle=\ttfamily\small,
  backgroundcolor=\color{cinza},
  frame=single,
  rulecolor=\color{black!20},
  columns=fullflexible,
  keepspaces=true,
  showstringspaces=false,
  numbers=left,
  numberstyle=\tiny\color{cinzaEscuro},
  numbersep=7pt,
  literate={á}{{\'a}}1 {ã}{{\~a}}1 {ç}{{\c c}}1 {é}{{\'e}}1
           {í}{{\'i}}1 {ó}{{\'o}}1 {ú}{{\'u}}1 {Á}{{\'A}}1
           {É}{{\'E}}1 {Í}{{\'I}}1 {Ó}{{\'O}}1 {Ú}{{\'U}}1
           {à}{{\`a}}1 {õ}{{\~o}}1 {â}{{\^a}}1 {ê}{{\^e}}1
           {ô}{{\^o}}1 {<=}{{$\le$}}1 {>=}{{$\ge$}}1 {!=}{{$\neq$}}1
           {<-}{{$\leftarrow$}}1
}

\begin{document}

\begin{center}
  {\small Universidade Federal de Uberlândia -- UFU\\Faculdade de Engenharia Elétrica}\par
  \vspace{4pt}
  {\LARGE\bfseries\color{azul} Lista 7: O Problema da Seleção}\par
  \vspace{4pt}
  \textbf{Arthur Pires Lopes}\quad Matrícula: \textbf{12221EEL001}\par
  \textit{Caderno Completo de Estudos: Demonstrações, Análises Detalhadas e Revisão Conceitual}
\end{center}
\hrule
\vspace{8pt}

% ==============================================================================
% GUIA DE ESTUDOS E MAPA CONCEITUAL
% ==============================================================================
\section*{Guia de Estudos e Fundamentação Teórica}

O \textbf{Problema da Seleção} consiste em encontrar a $i$-ésima menor chave em um conjunto não ordenado de $n$ elementos distintos (a chamada $i$-\textit{ésima estatística de ordem}).
\begin{itemize}
    \item Para $i = 1$: temos o \textbf{Mínimo};
    \item Para $i = n$: temos o \textbf{Máximo};
    \item Para $i = \lfloor (n+1)/2 \rfloor$ ou $i = \lceil (n+1)/2 \rceil$: temos a \textbf{Mediana}.
\end{itemize}

\begin{caixaEstudo}[Comparação Fundamental: Ordenação vs. Seleção]
Ordenar o conjunto inteiro consome $\Theta(n \lg n)$ no modelo de comparações. Contudo, para encontrar um único elemento (ou um pequeno grupo de elementos), ordenar tudo é redundante: a teoria garante que o problema da seleção pode ser resolvido em \textbf{tempo linear $\Theta(n)$ no pior caso}.
\end{caixaEstudo}

\topico{Diferenças entre Divisão e Conquista nos Algoritmos Clássicos}
\begin{center}
\small
\setlength{\tabcolsep}{4pt}
\begin{tabular}{p{2.4cm}p{2.8cm}p{3.8cm}p{5.8cm}}
\toprule
\textbf{Algoritmo} & \textbf{Paradigma} & \textbf{Complexidade} & \textbf{Comportamento da Recursão} \\ \midrule
\texttt{QuickSort} & Divisão e Conquista & $\Theta(n \lg n)$ médio, $O(n^2)$ pior & Conquista \textbf{ambos} os subproblemas ($2$ chamadas). \\
\texttt{Random-Select} & Podar e Buscar & $\Theta(n)$ médio, $O(n^2)$ pior & Descarta um lado e busca em \textbf{apenas um} subproblema. \\
\texttt{BFPRT (Select)} & Podar e Buscar Ótimo & $\Theta(n)$ pior caso & Mediana das medianas garante descarte de $\ge 30\%$ das chaves. \\
\bottomrule
\end{tabular}
\end{center}

\vspace{10pt}
\hrule
\vspace{6pt}

% ==============================================================================
% QUESTÃO 1
% ==============================================================================
\questao{1}{Mínimo e Máximo Simultâneos com no Máximo $3\lfloor n/2 \rfloor$ Comparações}

\textbf{Enunciado:} Escreva o pseudocódigo de algoritmo que encontre o máximo e o mínimo em um vetor $A[1, n]$ realizando no máximo $3\lfloor n/2 \rfloor$ comparações.

\topico{1. Motivação e Abordagem Ingênua}
Se buscarmos o mínimo e o máximo de maneira independente:
\begin{enumerate}
    \item Determinar o mínimo percorrendo o vetor: $n - 1$ comparações.
    \item Determinar o máximo percorrendo o vetor: $n - 1$ comparações.
\end{enumerate}
Total ingênuo: $2n - 2$ comparações. Para $n = 1000$, seriam $1998$ comparações.

\topico{2. Estratégia dos Pares (Torneio Parcial)}
Podemos economizar comparações mantendo um $\min$ e um $\max$ correntes e processando os elementos restantes \textbf{dois a dois} (em pares):
\begin{itemize}
    \item Seja o par $(A[i], A[i+1])$. Comparamos $A[i]$ com $A[i+1]$ (\destaque{1ª comparação}).
    \item O elemento que perdeu essa comparação (o menor do par) não pode ser o máximo do vetor. Portanto, comparamos ele unicamente com o $\min$ corrente (\destaque{2ª comparação}).
    \item O elemento que venceu essa comparação (o maior do par) não pode ser o mínimo do vetor. Logo, comparamos ele unicamente com o $\max$ corrente (\destaque{3ª comparação}).
\end{itemize}
\textbf{Conclusão teórica:} Para cada $2$ elementos processados, realizamos exatamente $3$ comparações (taxa de $1{,}5$ comparações por elemento, contra $2$ da ingênua).

\topico{3. Pseudocódigo Completo}
\begin{lstlisting}
MIN-MAX-OTIMO(A, n)
 1  se n == 1 então
 2      retorne (A[1], A[1])
 3  // Inicialização baseada na paridade de n
 4  se n é ímpar então
 5      min <- A[1]
 6      max <- A[1]
 7      inicio <- 2
 8  senão
 9      se A[1] <= A[2] então
10          min <- A[1]; max <- A[2]
11      senão
12          min <- A[2]; max <- A[1]
13      inicio <- 3
14  // Laço principal: processamento em pares
15  para i <- inicio até n com passo 2 faça
16      se A[i] <= A[i+1] então
17          se A[i] < min então min <- A[i]
18          se A[i+1] > max então max <- A[i+1]
19      senão
20          se A[i+1] < min então min <- A[i+1]
21          se A[i] > max então max <- A[i]
22  retorne (min, max)
\end{lstlisting}

\topico{4. Dedução Matemática da Quantidade de Comparações $C(n)$}
Analisamos rigorosamente os dois casos de paridade:

\textbf{Caso 1: $n$ é ímpar}
\begin{itemize}
    \item Na inicialização: nenhuma comparação entre chaves é realizada (linhas 5 e 6 apenas atribuem $A[1]$ a $\min$ e $\max$).
    \item Número de elementos restantes: $n - 1$ elementos (quantidade par).
    \item Quantidade de pares processados no laço: $\frac{n - 1}{2} = \lfloor n/2 \rfloor$.
    \item Cada par consome exatamente 3 comparações (linhas 16, 17/18 ou 16, 20/21).
    \item Custo total:
    \[
    C(n) = 0 + 3 \cdot \left(\frac{n - 1}{2}\right) = 3 \left\lfloor \frac{n}{2} \right\rfloor.
    \]
\end{itemize}

\textbf{Caso 2: $n$ é par}
\begin{itemize}
    \item Na inicialização: $A[1]$ e $A[2]$ são comparados entre si (linha 9), consumindo \destaque{1 comparação}.
    \item Número de elementos restantes: $n - 2$ elementos.
    \item Quantidade de pares no laço: $\frac{n - 2}{2} = \frac{n}{2} - 1 = \lfloor n/2 \rfloor - 1$.
    \item Custo das comparações no laço: $3 \times \left(\frac{n}{2} - 1\right)$.
    \item Custo total:
    \[
    C(n) = 1 + 3\left(\frac{n}{2} - 1\right) = 1 + \frac{3n}{2} - 3 = \frac{3n}{2} - 2 = 3\left\lfloor \frac{n}{2} \right\rfloor - 2.
    \]
\end{itemize}

\begin{caixaEstudo}[Resultado Final e Limite Superior]
Como para $n$ ímpar temos $C(n) = 3\lfloor n/2 \rfloor$ e para $n$ par temos $C(n) = 3\lfloor n/2 \rfloor - 2 \le 3\lfloor n/2 \rfloor$, concluímos:
\[
\boxed{C(n) \le 3\left\lfloor \frac{n}{2} \right\rfloor \quad \forall\, n \ge 1.}
\]
\textbf{Nota Teórica:} Um argumento de adversário em Teoria da Complexidade demonstra que qualquer algoritmo baseado em comparações precisa de pelo menos $\lceil 3n/2 \rceil - 2$ comparações no pior caso. Logo, este algoritmo atinge o \textbf{limite inferior ótimo da teoria da informação}!
\end{caixaEstudo}

\vspace{10pt}
\hrule
\vspace{6pt}

% ==============================================================================
% QUESTÃO 2
% ==============================================================================
\questao{2}{Execução Passo a Passo do Algoritmo \texttt{SELECTION}}

\textbf{Enunciado:} Mostre a execução passo-a-passo do algoritmo \texttt{SELECTION} para encontrar o $3^\circ$ menor elemento para o seguinte vetor:
\[
A = \langle 1, 8, 7, 6, 5, 3, 2, 10, 9 \rangle
\]
Considere que o pivot é escolhido sempre como o primeiro elemento do vetor.

\topico{1. Mecânica da Partição com Pivô no Início (\texttt{PARTITION-PRIMEIRO})}
Para entender a movimentação exata do vetor, analisamos o algoritmo de partição de Lomuto parametrizado para o primeiro elemento $A[p]$:
\begin{lstlisting}
PARTITION-PRIMEIRO(A, p, r)
 1  x <- A[p]         // pivo escolhido como o primeiro elemento
 2  m <- p            // indice limite da regiao dos elementos <= x
 3  para j <- p + 1 até r faça
 4      se A[j] <= x então
 5          m <- m + 1
 6          troca A[m] com A[j]
 7  troca A[p] com A[m] // posiciona o pivo entre as regioes
 8  retorne m          // posicao final definitiva do pivo
\end{lstlisting}

\topico{2. Invariante de Laço da Partição}
No início de cada iteração do laço \texttt{para} (linha 3):
\begin{itemize}
    \item $A[p] = x$ é o pivô;
    \item Todo elemento em $A[p+1 \dots m]$ é $\le x$;
    \item Todo elemento em $A[m+1 \dots j-1]$ é $> x$;
    \item Os elementos em $A[j \dots r]$ ainda não foram examinados.
\end{itemize}

\topico{3. Rastreamento Minucioso das 5 Chamadas}
Vetor original: $A = \langle 1, 8, 7, 6, 5, 3, 2, 10, 9 \rangle$, $n = 9$, buscamos a $3^{\underline{a}}$ menor estatística ($i = 3$).

\vspace{4pt}
\destaque{Chamada 1: \texttt{SELECTION(A, 1, 9, 3)}}
\begin{itemize}
    \item Intervalo ativo: $p = 1, r = 9$. Buscamos $i = 3$.
    \item Pivô selecionado: $x = A[1] = 1$. Ponteiro de menores: $m = 1$.
    \item Varredura de $j = 2$ a $9$:
    \[
    A[2]=8 > 1, \; A[3]=7 > 1, \; \dots, \; A[9]=9 > 1.
    \]
    Nenhum elemento é menor ou igual a $1$. Nenhuma troca é realizada no laço ($m$ permanece $1$).
    \item Troca final: $A[p] \leftrightarrow A[m] \implies A[1] \leftrightarrow A[1]$. O pivô permanece na posição $q = 1$.
    \item Vetor após partição: $[\;\underline{1}\;, \; 8, 7, 6, 5, 3, 2, 10, 9\;]$.
    \item Posto relativo do pivô: $k = q - p + 1 = 1 - 1 + 1 = 1$.
    \item Tomada de decisão: Como $i = 3 > k = 1$, o $3^\circ$ menor elemento é estritamente maior que $1$ e encontra-se na partição direita.
    \item Nova ordem relativa a buscar: $i' = i - k = 3 - 1 = 2$.
    \item Chamada recursiva: \texttt{SELECTION(A, 2, 9, 2)}.
\end{itemize}

\vspace{4pt}
\destaque{Chamada 2: \texttt{SELECTION(A, 2, 9, 2)}}
\begin{itemize}
    \item Intervalo ativo: $p = 2, r = 9$. Buscamos o $2^\circ$ menor deste subvetor ($i = 2$).
    \item Subvetor atual: $A[2 \dots 9] = \langle 8, 7, 6, 5, 3, 2, 10, 9 \rangle$.
    \item Pivô: $x = A[2] = 8$. Inicializa $m = 2$.
    \item Varredura de $j = 3$ a $9$:
    \begin{itemize}
        \item $j=3$: $A[3]=7 \le 8 \implies m=3$, troca $A[3] \leftrightarrow A[3]$ (mantém 7).
        \item $j=4$: $A[4]=6 \le 8 \implies m=4$, troca $A[4] \leftrightarrow A[4]$ (mantém 6).
        \item $j=5$: $A[5]=5 \le 8 \implies m=5$, troca $A[5] \leftrightarrow A[5]$ (mantém 5).
        \item $j=6$: $A[6]=3 \le 8 \implies m=6$, troca $A[6] \leftrightarrow A[6]$ (mantém 3).
        \item $j=7$: $A[7]=2 \le 8 \implies m=7$, troca $A[7] \leftrightarrow A[7]$ (mantém 2).
        \item $j=8$: $A[8]=10 > 8 \implies$ não entra na condição, $m$ permanece 7.
        \item $j=9$: $A[9]=9 > 8 \implies$ não entra na condição, $m$ permanece 7.
    \end{itemize}
    \item Troca final: $A[p] \leftrightarrow A[m] \implies A[2]$ troca com $A[7]$. O elemento $8$ vai para a posição $7$ e o elemento $2$ vai para a posição $2$!
    \item Vetor global atualizado:
    \[
    A = \langle 1, \;\mathbf{2, 7, 6, 5, 3},\; \underline{8},\; \mathbf{10, 9} \rangle.
    \]
    \item Posição final do pivô: $q = 7$.
    \item Posto relativo no subvetor: $k = q - p + 1 = 7 - 2 + 1 = 6$ (o $8$ é o $6^\circ$ menor de $A[2 \dots 9]$).
    \item Tomada de decisão: Buscamos $i = 2$. Como $i = 2 < k = 6$, o elemento desejado está à esquerda do pivô.
    \item Ordem a buscar: $i = 2$ (mantém o mesmo valor na recursão à esquerda).
    \item Chamada recursiva: \texttt{SELECTION(A, 2, 6, 2)}.
\end{itemize}

\vspace{4pt}
\destaque{Chamada 3: \texttt{SELECTION(A, 2, 6, 2)}}
\begin{itemize}
    \item Intervalo ativo: $p = 2, r = 6$. Buscamos $i = 2$.
    \item Subvetor atual: $A[2 \dots 6] = \langle 2, 7, 6, 5, 3 \rangle$.
    \item Pivô: $x = A[2] = 2$. Inicializa $m = 2$.
    \item Varredura de $j = 3$ a $6$:
    \[
    A[3]=7 > 2, \; A[4]=6 > 2, \; A[5]=5 > 2, \; A[6]=3 > 2.
    \]
    Nenhum elemento é $\le 2$. Nenhuma troca no laço; $m = 2$.
    \item Troca final: $A[2] \leftrightarrow A[2]$. Posição final: $q = 2$.
    \item Vetor global atualizado: $\langle 1, \;\underline{2},\; \mathbf{7, 6, 5, 3},\; 8, 10, 9 \rangle$.
    \item Posto relativo: $k = q - p + 1 = 2 - 2 + 1 = 1$.
    \item Tomada de decisão: Buscamos $i = 2$. Como $i = 2 > k = 1$, o elemento está à direita de $2$.
    \item Nova ordem a buscar: $i' = i - k = 2 - 1 = 1$.
    \item Chamada recursiva: \texttt{SELECTION(A, 3, 6, 1)}.
\end{itemize}

\vspace{4pt}
\destaque{Chamada 4: \texttt{SELECTION(A, 3, 6, 1)}}
\begin{itemize}
    \item Intervalo ativo: $p = 3, r = 6$. Buscamos o $1^\circ$ menor de $A[3 \dots 6]$ ($i = 1$).
    \item Subvetor atual: $A[3 \dots 6] = \langle 7, 6, 5, 3 \rangle$.
    \item Pivô: $x = A[3] = 7$. Inicializa $m = 3$.
    \item Varredura de $j = 4$ a $6$:
    \begin{itemize}
        \item $j=4$: $A[4]=6 \le 7 \implies m=4$, troca $A[4] \leftrightarrow A[4]$.
        \item $j=5$: $A[5]=5 \le 7 \implies m=5$, troca $A[5] \leftrightarrow A[5]$.
        \item $j=6$: $A[6]=3 \le 7 \implies m=6$, troca $A[6] \leftrightarrow A[6]$.
    \end{itemize}
    \item Troca final: $A[p] \leftrightarrow A[m] \implies A[3]$ troca com $A[6]$ (troca $7$ com $3$).
    \item Vetor global atualizado:
    \[
    A = \langle 1, 2, \;\mathbf{3, 6, 5},\; \underline{7},\; 8, 10, 9 \rangle.
    \]
    \item Posição final do pivô: $q = 6$.
    \item Posto relativo: $k = q - p + 1 = 6 - 3 + 1 = 4$.
    \item Tomada de decisão: Buscamos $i = 1$. Como $i = 1 < k = 4$, o elemento está à esquerda do pivô.
    \item Ordem a buscar: $i = 1$.
    \item Chamada recursiva: \texttt{SELECTION(A, 3, 5, 1)}.
\end{itemize}

\vspace{4pt}
\destaque{Chamada 5: \texttt{SELECTION(A, 3, 5, 1)}}
\begin{itemize}
    \item Intervalo ativo: $p = 3, r = 5$. Buscamos $i = 1$.
    \item Subvetor atual: $A[3 \dots 5] = \langle 3, 6, 5 \rangle$.
    \item Pivô: $x = A[3] = 3$. Inicializa $m = 3$.
    \item Varredura de $j = 4$ a $5$: $A[4]=6 > 3$ e $A[5]=5 > 3$. Nenhuma troca; $m = 3$.
    \item Troca final: $A[3] \leftrightarrow A[3]$. Posição final: $q = 3$.
    \item Vetor global atualizado: $\langle 1, 2, \;\underline{3},\; 6, 5, 7, 8, 10, 9 \rangle$.
    \item Posto relativo: $k = q - p + 1 = 3 - 3 + 1 = 1$.
    \item Tomada de decisão: Buscamos $i = 1$. Como $i == k = 1$, \destaque{o pivô é exatamente a resposta procurada!}
    \item O algoritmo retorna o elemento $A[q] = A[3] = 3$.
\end{itemize}

\topico{4. Tabela Síntese da Execução}
\begin{center}
\small
\setlength{\tabcolsep}{4pt}
\begin{tabular}{cclccccp{5.5cm}}
\toprule
\textbf{Passo} & \textbf{Intervalo} & \textbf{Vetor no Início} & \textbf{Pivô} & $q$ & $k$ & $i$ & \textbf{Ação Decidida} \\ \midrule
1 & $A[1 \dots 9]$ & $\langle 1, 8, 7, 6, 5, 3, 2, 10, 9 \rangle$ & $1$ & 1 & 1 & 3 & $i > k \implies$ Busca direita com $i'=2$ \\
2 & $A[2 \dots 9]$ & $\langle 8, 7, 6, 5, 3, 2, 10, 9 \rangle$ & $8$ & 7 & 6 & 2 & $i < k \implies$ Busca esquerda com $i=2$ \\
3 & $A[2 \dots 6]$ & $\langle 2, 7, 6, 5, 3 \rangle$ & $2$ & 2 & 1 & 2 & $i > k \implies$ Busca direita com $i'=1$ \\
4 & $A[3 \dots 6]$ & $\langle 7, 6, 5, 3 \rangle$ & $7$ & 6 & 4 & 1 & $i < k \implies$ Busca esquerda com $i=1$ \\
5 & $A[3 \dots 5]$ & $\langle 3, 6, 5 \rangle$ & $3$ & 3 & 1 & 1 & $i == k \implies$ \textbf{Retorna 3} \\
\bottomrule
\end{tabular}
\end{center}

\begin{caixaEstudo}[Verificação de Consistência]
Ordenando o vetor original: $\langle 1, 2, \mathbf{3}, 5, 6, 7, 8, 9, 10 \rangle$.
\begin{itemize}
    \item 1º menor: $1$;
    \item 2º menor: $2$;
    \item \destaque{3º menor: 3}.
\end{itemize}
O valor retornado pelo rastreamento é comprovadamente correto!
\end{caixaEstudo}

\vspace{10pt}
\hrule
\vspace{6pt}

% ==============================================================================
% QUESTÃO 3
% ==============================================================================
\questao{3}{Comparação dos Três Métodos para Listar os $i$ Maiores Elementos}

\textbf{Enunciado:} Dado um vetor $A[1, n]$, queremos imprimir em ordem crescente os $i$ maiores elementos de $A$. Compare a complexidade dos seguinte métodos em função de $n$ e $i$.
\begin{itemize}
    \item[(a)] Ordene o vetor e liste os $i$ maiores elementos.
    \item[(b)] Construa uma fila de prioridade (max-heap) e chame a rotina para extrair o máximo $i$ vezes.
    \item[(c)] Use um algoritmo de seleção para encontrar o $i$-ésimo maior elemento, particione o vetor em torno dele e ordene os $i$ maiores elementos.
\end{itemize}

\topico{1. Análise Formal do Método (a): Ordenação Global}
\begin{enumerate}
    \item \textbf{Ordenação de todo o vetor $A[1 \dots n]$:} Aplicando \texttt{MergeSort} ou \texttt{HeapSort}, o custo é $\Theta(n \lg n)$ tanto no pior caso quanto no caso médio.
    \item \textbf{Varredura dos $i$ maiores:} Os $i$ maiores elementos ocupam as posições $A[n - i + 1 \dots n]$. Como o vetor já está ordenado de forma crescente, basta percorrê-las e imprimi-las diretamente em tempo $\Theta(i)$.
\end{enumerate}
\[
T_a(n, i) = \Theta(n \lg n) + \Theta(i) = \Theta(n \lg n).
\]
\textit{Observação:} Como $i \le n$, o termo $n \lg n$ domina estritamente o termo $i$.

\topico{2. Análise Formal do Método (b): Max-Heap e $i$ Extrações}
\begin{enumerate}
    \item \textbf{Construção do heap (\texttt{BUILD-MAX-HEAP}):} Utilizando o algoritmo linear clássico de Floyd (de baixo para cima), o tempo é comprovadamente linear:
    \[
    \sum_{h=0}^{\lfloor \lg n \rfloor} \left\lceil \frac{n}{2^{h+1}} \right\rceil O(h) = O\left(n \sum_{h=0}^\infty \frac{h}{2^h}\right) = \Theta(n).
    \]
    \item \textbf{Extrações sucessivas ($i$ chamadas a \texttt{EXTRACT-MAX}):}
    Cada extração retira o maior elemento atual na raiz e chama \texttt{MAX-HEAPIFY}, cujo custo é proporcional à altura da árvore, isto é, $O(\lg |H|)$. Para $i$ extrações:
    \[
    \sum_{j=0}^{i-1} O(\lg(n - j)) = \Theta(i \lg n).
    \]
    \item \textbf{Ordenação final da saída:} As extrações fornecem os $i$ elementos em ordem \textbf{decrescente}. Como o enunciado exige explicitamente a saída em \textbf{ordem crescente}, devemos armazenar esses elementos em um vetor auxiliar de tamanho $i$ e percorrê-lo de trás para frente (ou invertê-lo), consumindo $\Theta(i)$.
\end{enumerate}
\[
T_b(n, i) = \Theta(n) + \Theta(i \lg n) + \Theta(i) = \Theta(n + i \lg n).
\]

\topico{3. Análise Formal do Método (c): Seleção Linear + Ordenação Parcial}
\begin{enumerate}
    \item \textbf{Identificação da estatística de ordem:} O $i$-ésimo maior elemento é o elemento de posto $k = n - i + 1$ em ordem crescente. Utilizando o algoritmo \texttt{SELECT} (mediana das medianas, Seção 9.3 de Cormen), determinamos este elemento em tempo garantido $\Theta(n)$ no pior caso.
    \item \textbf{Particionamento:} Ao particionar o vetor em torno desta chave, os $i - 1$ elementos maiores ficam agrupados no subvetor final $A[n - i + 2 \dots n]$, e a própria chave na posição $n - i + 1$. Custo da partição: $\Theta(n)$.
    \item \textbf{Ordenação restrita aos $i$ elementos:} Em vez de ordenar todo o vetor de tamanho $n$, ordenamos exclusivamente o subvetor de tamanho $i$ com \texttt{MergeSort}: custo $\Theta(i \lg i)$.
    \item \textbf{Listagem:} Percorrer os $i$ elementos ordenados: $\Theta(i)$.
\end{enumerate}
\[
T_c(n, i) = \Theta(n) + \Theta(n) + \Theta(i \lg i) + \Theta(i) = \Theta(n + i \lg i).
\]

\topico{4. Comparação em Diferentes Regimes Assintóticos de $i$}
\begin{center}
\footnotesize
\setlength{\tabcolsep}{3.5pt}
\begin{tabular}{p{2.3cm}p{2.2cm}p{3.2cm}p{3.5cm}p{3.4cm}}
\toprule
\textbf{Regime de $i$} & \textbf{Método (a)} & \textbf{Método (b)} & \textbf{Método (c)} & \textbf{Melhor Escolha} \\ \midrule
$i = O(1)$ & $\Theta(n \lg n)$ & $\Theta(n)$ & $\Theta(n)$ & (b) e (c) empatam em $\Theta(n)$ \\
$i = \Theta(\lg n)$ & $\Theta(n \lg n)$ & $\Theta(n + \lg^2 n) = \Theta(n)$ & $\Theta(n + \lg n \lg \lg n) = \Theta(n)$ & (b) e (c) operam em $\Theta(n)$ \\
$i = \Theta(\sqrt{n})$ & $\Theta(n \lg n)$ & $\Theta(n)$ & $\Theta(n)$ & (b) e (c) operam em $\Theta(n)$ \\
$i = \Theta(n / \lg n)$ & $\Theta(n \lg n)$ & $\Theta(n)$ & $\Theta(n)$ & (b) e (c) operam em $\Theta(n)$ \\
$i = \Theta(n)$ & $\Theta(n \lg n)$ & $\Theta(n \lg n)$ & $\Theta(n \lg n)$ & Empate assintótico \\
\bottomrule
\end{tabular}
\end{center}

\begin{caixaEstudo}[Conclusão Comparativa]
Como $1 \le i \le n$, temos sempre $\lg i \le \lg n$. Portanto, a relação de dominância é:
\[
\boxed{\Theta(n + i \lg i) \le \Theta(n + i \lg n) \le \Theta(n \lg n).}
\]
\textbf{Síntese Prática:}
\begin{itemize}
    \item O \textbf{Método (c)} é a melhor escolha assintótica teórica global: nunca perde para os outros e é estritamente superior para $i$ intermediário.
    \item O \textbf{Método (b)} é muito simples de implementar e extremamente veloz na prática quando $i$ é pequeno ($i \ll n$), pois o heap opera com constantes baixas.
    \item O \textbf{Método (a)} é ingênuo e subótimo sempre que $i = o(n)$.
\end{itemize}
\end{caixaEstudo}

\vspace{10pt}
\hrule
\vspace{6pt}

% ==============================================================================
% QUESTÃO 4
% ==============================================================================
\questao{4}{QuickSort Modificado com Tempo $\Theta(n \lg n)$ no Pior Caso}

\textbf{Enunciado:} Mostre como modificar \texttt{QuickSort} de modo que tenha complexidade de tempo $\Theta(n \lg n)$ no pior caso.

\topico{1. A Falha Estrutural do QuickSort Convencional}
No \texttt{QuickSort} padrão, se o pivô for sistematicamente escolhido como a menor ou maior chave restante (por exemplo, vetor já ordenado ou em ordem inversa com pivô no início ou no fim):
\begin{itemize}
    \item O particionamento divide o vetor de tamanho $m$ em subproblemas de tamanhos $0$ e $m - 1$.
    \item A relação de recorrência resultante é:
    \[
    T(n) = T(n-1) + T(0) + \Theta(n) = T(n-1) + \Theta(n).
    \]
    \item Pela soma dos primeiros inteiros, $T(n) = \sum_{k=1}^n \Theta(k) = \Theta(n^2)$.
\end{itemize}

\topico{2. O Mecanismo da Modificação: Seleção Determinística do Pivô}
Para garantir tempo ótimo no pior caso, devemos impor que a partição seja sempre balanceada. A chave para isso é usar o algoritmo determinístico \texttt{SELECT} (Cormen Seção 9.3) para encontrar a \destaque{mediana exata} do subvetor antes de chamar o particionamento.

\topico{3. Pseudocódigo Completo}
\begin{lstlisting}
QUICKSORT-PIOR-CASO-OTIMO(A, p, r)
 1  se p < r então
 2      m <- r - p + 1
 3      k <- piso((m + 1) / 2)               // posto da mediana exata
 4      mediana <- SELECT(A, p, r, k)        // algoritmo BFPRT: O(m) no pior caso
 5      localiza indice idx em [p, r] onde A[idx] == mediana
 6      troca A[idx] com A[r]               // posiciona a mediana no final
 7      q <- PARTITION(A, p, r)              // particiona em torno da mediana
 8      QUICKSORT-PIOR-CASO-OTIMO(A, p, q - 1)
 9      QUICKSORT-PIOR-CASO-OTIMO(A, q + 1, r)
\end{lstlisting}

\topico{4. Demonstração Formal da Complexidade}
\begin{enumerate}
    \item \textbf{Balanceamento Perfeito:} Como o pivô $A[q]$ é a mediana exata de $A[p \dots r]$:
    \begin{itemize}
        \item O subvetor esquerdo $A[p \dots q-1]$ contém $\lfloor (m - 1)/2 \rfloor \le m/2$ elementos.
        \item O subvetor direito $A[q+1 \dots r]$ contém $\lceil (m - 1)/2 \rceil \le m/2$ elementos.
    \end{itemize}
    \item \textbf{Custo Não Recursivo Linear:}
    \begin{itemize}
        \item A chamada \texttt{SELECT(A, p, r, k)} consome $c_1 m$ operações no pior caso.
        \item A localização do índice e a troca consomem $c_2 m$.
        \item O particionamento consome $c_3 m$.
        \item Total do nível: $c \cdot m = \Theta(m)$.
    \end{itemize}
    \item \textbf{Recorrência do Pior Caso:}
    \[
    T(n) \le 2\,T\left(\frac{n}{2}\right) + c\,n.
    \]
\end{enumerate}

\topico{5. Resolução da Recorrência pelo Teorema Mestre}
A recorrência é da forma $T(n) = a T(n/b) + f(n)$ com $a = 2$, $b = 2$ e $f(n) = cn$.
\begin{itemize}
    \item Expoente crítico: $n^{\log_b a} = n^{\log_2 2} = n^1 = n$.
    \item Comparação com $f(n)$: Como $f(n) = cn \in \Theta(n) = \Theta(n^{\log_b a})$, a recorrência se enquadra precisamente no \textbf{Caso 2 do Teorema Mestre}.
    \item Solução:
    \[
    T(n) \in \Theta(n^{\log_b a} \lg n) = \Theta(n \lg n).
    \]
\end{itemize}

\begin{caixaEstudo}[Por que este algoritmo não é usado nas bibliotecas padrão?]
Embora assintoticamente brilhante ($\Theta(n \lg n)$ no pior caso), o algoritmo \texttt{SELECT} possui uma constante oculta grande devido à divisão em blocos de 5 e ao overhead recursivo. Na prática da engenharia de software (como no \texttt{std::sort} de C++), utiliza-se o algoritmo híbrido \textbf{Introsort} (Musser, 1997):
\begin{enumerate}
    \item Executa o \texttt{QuickSort} padrão (com pivô mediano de 3);
    \item Monitora a profundidade da recursão;
    \item Se a profundidade ultrapassar $2\lfloor \lg n \rfloor$ (sinal de partição degenerada), aborta o QuickSort e comuta para o \texttt{HeapSort}.
\end{enumerate}
Dessa forma, o Introsort obtém o melhor dos dois mundos: a velocidade prática do QuickSort com a garantia de $\Theta(n \lg n)$ no pior caso do HeapSort.
\end{caixaEstudo}

\vspace{10pt}
\hrule
\vspace{6pt}

% ==============================================================================
% QUESTÃO 5
% ==============================================================================
\questao{5}{Algoritmo Linear de Seleção com Caixa-Preta de Mediana}

\textbf{Enunciado:} Suponha que você tenha uma subrotina do tipo ``caixa-preta'' que determina a mediana em tempo linear (no pior caso). Descreva um algoritmo linear simples que resolve o problema da seleção para todo $i$.

\topico{1. Paradigma \textit{Prune-and-Search} (Podar e Buscar)}
Diferente da ordenação (que precisa processar ambos os lados da partição), a busca por um posto específico só precisa continuar em \textbf{um único} lado da partição. O outro lado é permanentemente descartado (podado).

\topico{2. Pseudocódigo Completo}
\begin{lstlisting}
SELECAO-CAIXA-PRETA(A, p, r, i)
 1  se p == r então retorne A[p]
 2  // Passo 1: Obter a mediana via sub-rotina de tempo linear
 3  med <- MEDIANA(A, p, r)
 4  // Passo 2: Localizar e trocar para usar como pivo
 5  localiza indice idx em [p, r] onde A[idx] == med
 6  troca A[idx] com A[r]
 7  // Passo 3: Particionar o vetor
 8  q <- PARTITION(A, p, r)
 9  k <- q - p + 1   // rank relativo da mediana no subvetor
10  // Passo 4: Podar e buscar recursivamente
11  se i == k então
12      retorne A[q]
13  senão se i < k então
14      retorne SELECAO-CAIXA-PRETA(A, p, q - 1, i)
15  senão
16      retorne SELECAO-CAIXA-PRETA(A, q + 1, r, i - k)
\end{lstlisting}

\topico{3. Dedução da Recorrência no Pior Caso}
Seja $m = r - p + 1$. Como o pivô é a mediana exata:
\begin{itemize}
    \item Pelo menos $\lfloor m/2 \rfloor$ elementos são menores ou iguais a ele.
    \item Pelo menos $\lfloor m/2 \rfloor$ elementos são maiores ou iguais a ele.
\end{itemize}
Portanto, qualquer que seja a metade escolhida (esquerda ou direita), o tamanho do novo subproblema é no máximo $\lfloor m/2 \rfloor \le m/2$.

Como o algoritmo realiza apenas \textbf{uma} chamada recursiva, e o trabalho não recursivo (caixa-preta + busca linear + partição) consome tempo $c \cdot n$, a relação de recorrência para o pior caso é:
\[
T(n) \le T\left(\frac{n}{2}\right) + c\,n.
\]

\topico{4. Prova 1: Método da Expansão e Soma Geométrica Infinita}
Expandindo a relação sucessivamente:
\begin{align*}
T(n) &\le T(n/2) + cn \\
     &\le \left[T(n/4) + c\left(\frac{n}{2}\right)\right] + cn \\
     &\le T(n/8) + c\left(\frac{n}{4}\right) + c\left(\frac{n}{2}\right) + cn \\
     &\le cn \left[ 1 + \frac{1}{2} + \frac{1}{4} + \frac{1}{8} + \dots + \left(\frac{1}{2}\right)^{\lfloor \lg n \rfloor} \right] + T(1).
\end{align*}
Como a série geométrica infinita converge:
\[
\sum_{j=0}^{\infty} \left(\frac{1}{2}\right)^j = \frac{1}{1 - 1/2} = 2.
\]
Portanto:
\[
T(n) \le cn \cdot 2 + O(1) = 2cn + O(1) = O(n).
\]
Como qualquer algoritmo de seleção precisa examinar todas as $n$ chaves no pior caso, $T(n) \in \Omega(n)$. Logo:
\res{T(n) \in \Theta(n).}

\topico{5. Prova 2: Aplicação do Teorema Mestre}
A recorrência é $T(n) = a T(n/b) + f(n)$ com $a = 1$, $b = 2$, e $f(n) = cn$.
\begin{itemize}
    \item Expoente crítico: $n^{\log_b a} = n^{\log_2 1} = n^0 = 1$.
    \item Como $f(n) = cn \in \Omega(n^{0 + \varepsilon})$ com $\varepsilon = 1$, verificamos a condição de regularidade:
    \[
    a\,f(n/b) = 1 \cdot c\left(\frac{n}{2}\right) = \frac{1}{2}\,cn \le d \cdot f(n), \quad \text{com } d = \frac{1}{2} < 1.
    \]
    A condição é satisfeita para todo $n \ge 1$.
    \item Pelo \textbf{Caso 3 do Teorema Mestre}:
    \[
    T(n) \in \Theta(f(n)) = \Theta(n).
    \]
\end{itemize}

\vspace{10pt}
\hrule
\vspace{6pt}

% ==============================================================================
% REVISÃO RÁPIDA / FLASHCARDS
% ==============================================================================
\section*{Revisão Rápida para a Prova (Perguntas e Respostas Chave)}

\begin{enumerate}
    \item \textbf{P: Por que o algoritmo de seleção ingênuo para achar min e max gasta $2n-2$ e o em pares gasta $3\lfloor n/2 \rfloor$?}
    
    \textit{R:} O algoritmo ingênuo compara cada elemento separadamente com min e max. O algoritmo em pares compara primeiro os dois elementos entre si, eliminando o perdedor da disputa do máximo e o vencedor da disputa do mínimo, gastando 3 comparações a cada 2 elementos.

    \item \textbf{P: No algoritmo \texttt{SELECTION}, por que ao recorrer para a direita usamos $i - k$?}
    
    \textit{R:} Porque ao descartar a partição esquerda e o pivô, já eliminamos exatamente $k$ elementos que são menores que todas as chaves da direita. Logo, a $i$-ésima menor chave global corresponde à $(i - k)$-ésima menor chave da região direita.

    \item \textbf{P: Qual é a intuição por trás da recorrência $T(n) = T(n/2) + O(n) = O(n)$?}
    
    \textit{R:} Como o trabalho não recursivo cai pela metade a cada passo ($n + n/2 + n/4 + \dots$), a soma é uma PG decrescente cuja soma total nunca ultrapassa o dobro do primeiro termo ($2n$). Logo, o primeiro passo domina todo o custo restante!

    \item \textbf{P: Por que o QuickSort pior caso com \texttt{SELECT} dá $\Theta(n \lg n)$ e a seleção com \texttt{SELECT} dá $\Theta(n)$?}
    
    \textit{R:} Porque o QuickSort possui $a = 2$ chamadas recursivas ($T(n) = 2T(n/2) + O(n)$), o que gera uma árvore com $\lg n$ níveis de custo constante $cn$, totalizando $cn \lg n$. Já a seleção possui $a = 1$ chamada recursiva ($T(n) = T(n/2) + O(n)$), de modo que o tamanho do problema se reduz geometricamente e a soma converge para $O(n)$.
\end{enumerate}

\end{document}
